\documentclass[10pt,a4paper]{article}
\usepackage[margin=1.7cm]{geometry}
\usepackage{enumitem}\usepackage{titlesec}\usepackage[hidelinks]{hyperref}\usepackage{parskip}
\titleformat{\section}{\normalfont\large\bfseries}{\thesection}{0.6em}{}
\titlespacing*{\section}{0pt}{6pt}{2pt}
\setlist[itemize]{topsep=2pt,itemsep=1pt,parsep=0pt,leftmargin=1.3em}
\pagestyle{empty}
\begin{document}
\begin{center}
{\Large\bfseries AVR Microcontroller Laboratory \textendash{} 10 Experiments in Embedded C}\\[2pt]
{\large Peripherals, interrupts and buses on the ATmega16/32 (CodeVisionAVR / Atmel Studio + Proteus)}\\[4pt]
{\small Digital Systems Design Laboratory 2 \textendash{} B.Sc. Electrical Engineering, Isfahan University of Technology}\\
{\small Author: Kiavash Rabbani Soudabi \textendash{} term 1401 (2022--2023)}
\end{center}

\section*{Purpose}
Ten laboratory experiments covering the peripherals of an 8-bit AVR microcontroller end to end: write the embedded C
driver, configure the peripheral at register level, simulate the circuit in Proteus and validate the behaviour on the
schematic. The experiments progress from GPIO up to interrupt-driven serial buses.

\section*{Experiments}
\begin{itemize}
  \item \textbf{1 \textendash{} Digital I/O ports:} reading switches and driving LEDs through the port registers.
  \item \textbf{2 \textendash{} Seven-segment displays:} multiplexed display driving with a modular source layout.
  \item \textbf{3 \textendash{} Alphanumeric LCD and keypad:} LCD initialisation and printing, 4x4 keypad scanning,
        and interaction between the two.
  \item \textbf{4 \textendash{} Timers and interrupts:} timer configuration, precise delays and interrupt service
        routines.
  \item \textbf{5 \textendash{} Motor drives:} DC and stepper motor control (direction, speed).
  \item \textbf{6 \textendash{} ADC:} the same conversion implemented both by polling and with interrupts, and the two
        approaches compared.
  \item \textbf{7 \textendash{} UART:} initialisation, interrupt-driven transmission/reception and printing received
        text on the LCD.
  \item \textbf{8 \textendash{} Graphic LCD:} 128x64 display with bitmap rendering, a clock and a dot-matrix font.
  \item \textbf{9 \textendash{} SPI:} master/slave communication between two microcontrollers with LCD output (group
        exercise).
  \item \textbf{10 \textendash{} I2C (TWI) and EEPROM:} master/slave I2C between an ATmega16 and an ATmega32, plus
        read/write of a 24LC32 EEPROM.
\end{itemize}

\section*{Verification and tools}
Embedded C written in CodeVisionAVR and Atmel Studio, simulated in Proteus with the in-circuit AVR model; each
experiment keeps its project files, its schematic (\texttt{.pdsprj}) and the compiled \texttt{.hex} so the simulation
runs as submitted. My written reports for experiments 1, 3 and 5 are included, together with screen captures of the
Proteus schematics and a compressed screen recording per experiment.

\section*{Repository contents}
One folder per experiment (\texttt{exp01-io-port} \dots \texttt{exp10-i2c-eeprom}) with the C sources, project files,
Proteus schematics, firmware images and demo videos; \texttt{docs/} with this overview.

\section*{Skills demonstrated}
Embedded C on an 8-bit AVR; register-level peripheral programming (GPIO, timers, ADC, UART, SPI, I2C/TWI);
interrupt-driven design and ISRs; display drivers for 7-segment, alphanumeric and graphic LCDs; motor control;
master/slave multi-microcontroller buses; EEPROM interfacing; the CodeVisionAVR and Atmel Studio toolchains;
simulation-based validation in Proteus.
\end{document}
